% !TeX root = paper.tex
% tables.tex —— 由 code/make_numbers.py 自动生成，请勿手工修改
\begin{table}[htbp]
\centering
\caption{调度中心与服务区基础数据（海拔取附件给定值，DEM 值为 30\,m 栅格最近邻采样）}
\label{tab:nodes}
\footnotesize
\resizebox{\linewidth}{!}{%
\begin{tabular}{llrrrrrr}
\toprule
编号 & 名称 & 经度/$^\circ$E & 纬度/$^\circ$N & 海拔/m & 人口 & DEM 高程/m & 直线距离/km \\
\midrule
O01 & 凤丹村无人机调度中心 & 109.2309 & 23.0085 & 127.7 & — & — & 0.00 \\
S001 & 平旺、那龙、镇龙乡 & 109.2432 & 23.0336 & 154.0 & 2100 & 154.0 & 3.06 \\
S002 & 上良 & 109.2586 & 23.0703 & 189.5 & 330 & 188.8 & 7.44 \\
S003 & 六造、那旭、都长 & 109.1689 & 23.0403 & 308.5 & 275 & 308.5 & 7.26 \\
S004 & 大站 & 109.2382 & 23.0778 & 198.3 & 188 & 195.1 & 7.75 \\
S005 & 上良、古律、那麓 & 109.2283 & 23.0561 & 165.0 & 126 & 169.0 & 5.30 \\
S006 & 六昌、木路 & 109.1998 & 23.0077 & 284.8 & 121 & 275.1 & 3.18 \\
S007 & 佛子、那桑、银坑 & 109.1928 & 23.0294 & 318.5 & 76 & 314.2 & 4.53 \\
S008 & 古麻、大站 & 109.2127 & 23.0795 & 240.9 & 55 & 248.3 & 8.11 \\
S009 & 那龙 & 109.2541 & 23.0553 & 200.0 & 45 & 189.4 & 5.72 \\
S010 & 镇龙乡 & 109.2760 & 23.0194 & 321.6 & 31 & 321.6 & 4.78 \\
S011 & 那从 & 109.2203 & 23.0319 & 182.1 & 30 & 178.0 & 2.82 \\
S012 & 镇龙乡 & 109.2831 & 23.0545 & 249.8 & 16 & 253.3 & 7.40 \\
S013 & 镇龙乡 & 109.2714 & 23.0311 & 321.7 & 14 & 321.7 & 4.85 \\
S014 & 镇龙乡 & 109.2836 & 23.0052 & 321.0 & 12 & 321.0 & 5.41 \\
S015 & 六造、那托、那旭 & 109.1924 & 23.0495 & 444.5 & 3 & 436.2 & 6.02 \\
\bottomrule
\end{tabular}
}
\end{table}

\begin{table}[htbp]
\centering
\caption{运输无人机 A/B/C 与中继无人机 R 的主要参数（$m_0$ 为含电池空载质量）}
\label{tab:types}
\footnotesize
\resizebox{\linewidth}{!}{%
\begin{tabular}{llrrrrrrrrrr}
\toprule
机型 & 名称 & $m_0$/kg & $Q$/kg & $V$/m$^3$ & $v_c$/(m$\cdot$s$^{-1}$) & $L^0$/km & $L^F$/km & $E^{\rm use}$/kWh & 架数 & 电池数 & $T_{\rm full}$/s \\
\midrule
A & 中轻载标准测试多旋翼 & 70.0 & 25 & 0.060 & 12.0 & 25.0 & 20.0 & 4.5 & 4 & 6 & 1800 \\
B & 中载标准测试多旋翼 & 65.0 & 30 & 0.073 & 15.0 & 28.0 & 16.0 & 4.0 & 2 & 4 & 2400 \\
C & 重载标准测试多旋翼 & 69.9 & 80 & 0.250 & 15.0 & 26.0 & 12.0 & 8.0 & 2 & 4 & 3000 \\
R & 中继标准测试多旋翼 & 23.5 & — & — & 15.0 & — & — & 3.2 & 2 & 6 & 1800 \\
\bottomrule
\end{tabular}
}
\end{table}

\begin{table}[htbp]
\centering
\caption{三机型在 15 个服务区的最大安全载荷（kg，$\rho=20\%$；$^{*}$ 表示受返航能量限制而非载重限制）}
\label{tab:payload}
\footnotesize
\begin{tabular}{lrrr}
\toprule
服务区 & A 型 & B 型 & C 型 \\
\midrule
S001 & 25.0 & 30.0 & 80.0 \\
S002 & 25.0 & 30.0 & 68.6$^{*}$ \\
S003 & 25.0 & 30.0 & 68.2$^{*}$ \\
S004 & 25.0 & 30.0 & 63.3$^{*}$ \\
S005 & 25.0 & 30.0 & 80.0 \\
S006 & 25.0 & 30.0 & 80.0 \\
S007 & 25.0 & 30.0 & 80.0 \\
S008 & 25.0 & 28.7$^{*}$ & 58.5$^{*}$ \\
S009 & 25.0 & 30.0 & 80.0 \\
S010 & 25.0 & 30.0 & 80.0 \\
S011 & 25.0 & 30.0 & 80.0 \\
S012 & 25.0 & 30.0 & 68.0$^{*}$ \\
S013 & 25.0 & 30.0 & 80.0 \\
S014 & 25.0 & 30.0 & 80.0 \\
S015 & 25.0 & 30.0 & 80.0 \\
\bottomrule
\end{tabular}
\end{table}

\begin{table}[htbp]
\centering
\caption{问题一五个具名组批方案的指标对比（累计作业时间为各架次时长之和）}
\label{tab:q1schemes}
\footnotesize
\resizebox{\linewidth}{!}{%
\begin{tabular}{lrrrrr}
\toprule
方案 & 架次数 & 总能耗/kWh & 累计工时/h & 平均质量占用率 & 平均体积占用率 \\
\midrule
S1-最少架次 & 18 & 59.23 & 7.99 & 0.787 & 0.772 \\
S4-最省能耗 & 19 & 59.13 & 8.48 & 0.795 & 0.797 \\
S5-最短总工时 & 18 & 59.23 & 7.99 & 0.787 & 0.772 \\
S2-最少架次·能耗最优 & 18 & 59.23 & 7.99 & 0.787 & 0.772 \\
S3-最少架次·工时最优 & 18 & 59.23 & 7.99 & 0.787 & 0.772 \\
\bottomrule
\end{tabular}
}
\end{table}

\begin{table}[htbp]
\centering
\caption{返航安全余量 $\rho$ 灵敏度：载荷上限与全场景架次数（「不可行」表示该 $\rho$ 下 S004/S008 连空载往返都不满足能量约束，全场景无可行组批方案）}
\label{tab:rho}
\footnotesize
\begin{tabular}{lrrrrrr}
\toprule
$\rho$ & A 型均值/kg & B 型均值/kg & C 型均值/kg & C 型受限区数 & 最少架次数 & 总能耗/kWh \\
\midrule
10\% & 25.0 & 30.0 & 78.7 & 5 & 18 & 59.23 \\
15\% & 25.0 & 30.0 & 77.1 & 5 & 18 & 59.23 \\
20\% & 25.0 & 29.9 & 75.1 & 5 & 18 & 59.23 \\
25\% & 24.9 & 29.4 & 72.4 & 6 & 19 & 61.18 \\
30\% & 23.9 & 28.1 & 68.5 & 7 & 20 & 67.34 \\
35\% & 19.9 & 26.4 & 62.0 & 9 & 25 & 78.16 \\
40\% & 16.7 & 23.2 & 52.7 & 11 & 不可行 & 不可行 \\
\bottomrule
\end{tabular}
\end{table}

\begin{table}[htbp]
\centering
\caption{问题二五种权重配置的 ALNS + CP-SAT 求解结果对比}
\label{tab:q2configs}
\footnotesize
\resizebox{\linewidth}{!}{%
\begin{tabular}{lrrrrrrrr}
\toprule
方案 & 权重$^*$ & 架次数 & 能耗/kWh & 完工时间/h & 延误量/h & 及时率 & 覆盖箱数 & 硬约束违反 \\
\midrule
P1-架次优先 & 4:1:1:1 & 17 & 59.87 & 2.64 & 70.04 & 91.8\% & 73 & 0 \\
P4-完工优先 & 0.5:0.5:6:1 & 18 & 67.76 & 2.84 & 56.32 & 86.2\% & 80 & 0 \\
P3-能耗优先 & 1:6:1:1 & 19 & 67.38 & 3.30 & 85.82 & 86.2\% & 80 & 0 \\
P5-均衡 & 1.5:2:2:2 & 19 & 68.16 & 2.88 & 44.38 & 82.5\% & 80 & 0 \\
P2-及时优先 & 0.5:1:0.5:6 & 20 & 73.31 & 3.80 & 30.94 & 81.2\% & 80 & 0 \\
\bottomrule
\end{tabular}
}
\\[2pt] \footnotesize $^*$ 权重顺序为「架次数 : 能耗 : 完工时间 : 加权延误」。五组权重下得到的均为**非支配解**，不构成全局 Pareto 前沿的证明。
\end{table}

\begin{table}[htbp]
\centering
\caption{问题三中继架次明细 —— \textbf{库存口径（2 架中继）下的最优折中方案}，共 3 个架次；达成全程连续通信所需的 3 架增配方案见 results/结果提交.xlsx 的「Q3\_中继架次\_增配方案」表}
\label{tab:q3relay}
\scriptsize
\resizebox{\linewidth}{!}{%
\begin{tabular}{lllrrrrrrrrr}
\toprule
编号 & 中继机 & 能源组件 & 悬停经度 & 悬停纬度 & 悬停海拔/m & 离地/m & 起飞/s & 建链/s & 服务结束/s & 返航/s & 能耗/kWh \\
\midrule
Q3-R-01 & R01 & R-E1 & 109.2639 & 23.0547 & 351 & 80 & 214 & 710 & 1970 & 2454 & 0.67 \\
Q3-R-02 & R02 & R-E2 & 109.2709 & 23.0057 & 455 & 80 & 2034 & 2420 & 3820 & 4203 & 0.63 \\
Q3-R-03 & R01 & R-E3 & 109.2359 & 23.0652 & 600 & 300 & 6030 & 6600 & 8800 & 9379 & 0.98 \\
\bottomrule
\end{tabular}
}
\end{table}

\begin{table}[htbp]\centering
\caption{问题一组批方案复算校验报告（results/Q1\_复算报告.txt）}
\label{tab:q1check}
\small
\begin{tabular}{l}
\toprule
架次数=18 覆盖货箱=80/80 违反项=0 \\
\bottomrule
\end{tabular}
\end{table}

\begin{table}[htbp]
\centering
\caption{论文主要符号说明}
\label{tab:symbols}
\small
\begin{tabular}{lp{11.5cm}}
\toprule
符号 & 含义 \\
\midrule
$G$ & 重力加速度，取 9.80665 m/s$^2$ \\
$q$ & 架次装载质量（kg） \\
$Q_g$ & 机型 $g$ 的最大载货质量（kg） \\
$L_g(q)$ & 机型 $g$ 在载荷 $q$ 下的等效航程（m） \\
$E_g^{\rm use}$ & 机型 $g$ 单组电池可用能量（kWh） \\
$E^{\rm hor}$ & 航段水平飞行能耗（kWh） \\
$E^{\rm up}$ & 航段爬升附加能耗（kWh） \\
$\eta_{\rm up}$ & 爬升能耗效率，取 0.72 \\
$h^{+},h^{-}$ & 航段爬升/下降高度（m） \\
$v_c,v_{\uparrow},v_{\downarrow}$ & 巡航/爬升/下降速度（m/s） \\
$\rho_g$ & 返航安全余量比例，三种运输机型均取 20\% \\
$P_{\rm th}$ & 接收门限（dBm），$P_{\rm th}=P_{\rm sens}+M=-90$ dBm \\
$L_{\max}^{a\to b}$ & 链路 $a\to b$ 的最大允许传播损耗（dB） \\
$L_{\rm obs}$ & 地形遮挡附加损耗（dB），取 10 dB \\
$\Delta$ & 航迹与链路的采样步长（s 或 m） \\
$x_{k,a}$ & 架次 $k$ 是否访问服务区 $a$ 的 0--1 变量 \\
\bottomrule
\end{tabular}
\end{table}

